\section{Part I: What Few Target Labels Fix First}
\label{sec:part1-what-few-labels-fix}

Part~I asks two questions about the linear head of Section~\ref{sec:setting}: what $K$ labels per emotion change in it, and how much of that change a source corpus can supply instead.
We first describe the gain map, then state the assumptions under which a threshold and a direction have different estimation costs, and then test four predictions that follow from these costs.

\subsection{The Gain Map}
\label{sec:part1-map}

Adding a source helps when target labels are scarce and hurts when they are plentiful (Fig.~\ref{fig:gainmap}).
For audio, the mean gain at $K=1$ is $\NUM{map-gain-a-k1}$ points, and its 95\% interval lies above zero for \NUM{map-gain-pos-a-k1} of the \NUM{map-pairs-a} ordered pairs.
The mean gain falls below zero by $K=25$ ($\NUM{map-gain-a-k25}$ points), and with the full pool the interval lies below zero for \NUM{map-gain-neg-a-kall} pairs.
All eight configurations follow the same course: the mean gain at $K=1$ ranges from $\NUM{map-gain-v-k1}$ (V) to $\NUM{map-gain-avt-k1}$ points (AVT), and at $K=25$ it is negative in every configuration.
The size of the gain is largely a property of the two corpora.
An additive model with one effect per target and one per source explains \NUM{map-share-a-lo}--\NUM{map-share-v-hi}\% of the variance of the gains and of the largest budget at which a source still helps, across A, A~(HuBERT), V, and AV.
The rest of Part~I explains this course through what the labels change in the classifier.

\begin{figure*}[t]
\centering
\FIG{gainmap}{\textwidth}
\caption{The gain map. (a)~Mean gain of pooling over training on the target labels alone (pooled minus target-only UAR in the same iteration), averaged over ordered pairs, against the label budget $K$ per emotion. (b)~Share of pairs whose 95\% interval of the gain lies above zero, includes zero, or lies below zero; these intervals treat the thirty iterations as independent, which makes them too narrow for IEMOCAP targets. All ordered pairs, including the two pairs of IEMOCAP sessions; the pooled classifier uses the penalty of the source-only classifier.}
\label{fig:gainmap}
\end{figure*}

\subsection{Two Costs of a Few Labels}
\label{sec:part1-costs}

We analyze the pairwise decisions of Section~\ref{sec:setting-model} under three assumptions.
\begin{enumerate}
\item[(A1)] Within each emotion, the score difference $z$ of a pair is Gaussian with a variance shared by the two emotions, so that \eqref{eq:twocloud} gives the pairwise accuracy.
\item[(A2)] The learner keeps a finite $\ell_2$ penalty on the weights and leaves the intercepts unpenalized. With a vanishing penalty on separable data, additive corrections of the logits lose their effect~\cite{CITE:kini2021}; a finite penalty keeps the intercepts free to move.
\item[(A3)] The $K$ labels per emotion are a random draw from the pool.
\end{enumerate}

\emph{Threshold.}
With the direction fixed, the intercept that balances two emotions places the threshold at the midpoint of the two class means of $z$.
Estimated from $K$ labels per emotion, this midpoint has variance $\sigma^2/(2K)$, so the residual offset satisfies $\mathbb{E}[o^2] = 1/(2K)$.
Expanding \eqref{eq:twocloud} to second order in $o$ gives the expected threshold loss that remains,
\begin{equation}
\label{eq:thrcost}
\mathbb{E}\left[a(d, 0) - a(d, o)\right] \approx \frac{d}{4}\,\phi\!\left(\frac{d}{2}\right)\frac{1}{2K} \le \frac{\phi(1)}{4K},
\end{equation}
where $\phi$ is the standard normal density and the maximum over $d$ is attained at $d=2$.
A handful of labels therefore removes most of any threshold offset, however large the offset was.

\emph{Direction.}
Estimating the direction requires estimating a vector, and its error shrinks with the ratio of an effective dimension $q$ to $K$.
For a linear rule estimated from $K$ labels per emotion, the separation behaves as
\begin{equation}
\label{eq:dircost}
d(K) = \frac{\Delta}{\sqrt{1 + q/K}},
\end{equation}
the form that the asymptotic expansion of the misclassification rate of the linear discriminant takes in its direction term~\cite{CITE:okamoto1963,CITE:raudys1998}.
We measure $q$ by fitting \eqref{eq:dircost} to the separations of target-only classifiers along the label budget.
The fitted $q$ is $\NUM{p1-qdir-a}$ for WavLM audio, $\NUM{p1-qdir-ah}$ for HuBERT audio, $\NUM{p1-qdir-av}$ for AV, and $\NUM{p1-qdir-v}$ for the face, so for audio a direction comparable to the target's best needs on the order of twenty labels per emotion.
In the ranges of $K$ studied here, labels thus fix thresholds cheaply and audio directions slowly.
Four predictions follow, and Table~\ref{tab:part1} collects the measurements that test them.
(P1)~Most of the loss of a source-only classifier that a few labels recover is threshold loss.
(P2)~The share of the gain that comes from the threshold falls as $K$ grows.
(P3)~Pooling corrects the threshold in proportion to the weight that the target receives (Section~\ref{sec:part1-pooling}).
(P4)~A source substitutes for a number of target labels set by $q$ and by how well its direction serves the target (Section~\ref{sec:part1-value}).

\begin{figure}[t]
\centering
\FIG{twocloud}{\columnwidth}
\caption{Threshold and direction of a pairwise decision. (a)~Two Gaussian classes along the score difference $z$; the threshold offset $o$ moves the decision point away from the midpoint, and the separation $d$ is set by the direction. (b)~Loss of source-only classifiers against the target's best linear rule, split into the threshold loss and the direction loss, in points of pairwise accuracy, averaged over the six emotion pairs and over ordered pairs.}
\label{fig:twocloud}
\end{figure}

\subsection{What the Labels Fix}
\label{sec:part1-fix}

The two-cloud model reproduces the pairwise accuracies of the classifiers we train.
Over all emotion pairs and ordered pairs, the accuracy predicted from the measured $d$ and $o$ correlates with the measured accuracy at \NUM{p1-e13-fit-r-lo}--\NUM{p1-e13-fit-r-hi}, with a mean absolute error of \NUM{p1-e13-fit-mae-lo}--\NUM{p1-e13-fit-mae-hi} points.

\emph{Source-only classifiers lose mainly through the direction, except on the face.}
For WavLM audio, the threshold loss is \NUM{p1-e13-loss-offset-a} points and the direction loss \NUM{p1-e13-loss-dir-a} points (Fig.~\ref{fig:twocloud}b); for the face the two are \NUM{p1-e13-loss-offset-v} and \NUM{p1-e13-loss-dir-v} points, and for text \NUM{p1-e13-loss-offset-t} and \NUM{p1-e13-loss-dir-t} points.
The threshold loss accounts for only \NUM{p1-e29-lossshare-lo}--\NUM{p1-e29-lossshare-hi} of the total loss across configurations.

\emph{The threshold loss is mostly a shift of the corpus mean.}
Centering, which subtracts the mean of the unlabeled utterances of each corpus and uses no label, raises the source-only UAR by \NUM{p1-center-srconly-a} points for WavLM audio, \NUM{p1-center-srconly-ah} for HuBERT audio, and \NUM{p1-center-srconly-v} for the face, and it raises it for \NUM{p1-center-srconly-up-a-n} of the \NUM{p1-center-srconly-pairs-a-n} audio pairs (three iterations).
Refitting only the intercepts of the source-only classifier on $K=5$ labels per emotion, with the weights frozen, gains \NUM{p1-d1-gain-a-k5} points for WavLM audio, close to the \NUM{p1-merged-gain-a-k5} points that pooling gains with the same labels.
After centering, the two gains shrink to \NUM{p1-center-d1-gain-a-k5} and \NUM{p1-center-merged-gain-a-k5} points (Fig.~S2 of the supplementary material).

\emph{The labels go to the threshold first (P1, P2).}
Writing the improvement of the pooled classifier over the source-only classifier as the sum of a threshold part and a direction part, computed from the measured $d$ and $o$ with \eqref{eq:twocloud}, the threshold part dominates for single-modality audio and face at $K \le 5$: its share is \NUM{p1-e29-critshare-kle5-single-lo}--\NUM{p1-e29-critshare-kle5-single-hi} across WavLM audio, HuBERT audio, and the face at $K=1$ and $K=5$ (Fig.~\ref{fig:thresholdshare}a and Table~\ref{tab:part1}; the shares exclude the two pairs of IEMOCAP sessions).
The share is much larger than the threshold's share of the loss, so the labels correct the cheap part of the error first.
As predicted, the share falls with $K$; at $K=10$ it is \NUM{p1-e29-critshare-k10-ah} for HuBERT audio.
In the fused and text configurations the threshold loss is smaller, and at $K=5$ the threshold part is about half of the gain or less (\NUM{p1-e29-critshare-k5-fused-text-lo}--\NUM{p1-e29-critshare-k5-fused-text-hi}, with \NUM{p1-e29-critshare-k5-av} for AV); at $K=1$ it still exceeds half for AV (\NUM{p1-e29-critshare-k1-av}) and VT (\NUM{p1-e29-critshare-k1-vt}).
A two-layer perceptron on the same features gives threshold shares of \NUM{p1-e30-critshare-k5-audio-lo}--\NUM{p1-e30-critshare-k5-audio-hi} for audio and \NUM{p1-e30-critshare-k5-v-wd2}--\NUM{p1-e30-critshare-k5-v-wd4} for the face at $K=5$, so the pattern is not specific to the linear head, although for HuBERT audio it falls below half.

\emph{The direction stays where the source left it.}
Under \eqref{eq:dircost}, even the optimal combination of the source direction with ten labels per emotion recovers \NUM{p1-dirloss-recov-bound-k10-audio-lo}--\NUM{p1-dirloss-recov-bound-k10-audio-hi}\% of the direction loss for audio, and hypothesis transfer, which starts from the source classifier and refits it on the labels, recovers \NUM{p1-dirloss-recov-ht-k10-audio-lo}--\NUM{p1-dirloss-recov-ht-k10-audio-hi}\%.

\begin{figure}[t]
\centering
\FIG{thresholdshare}{\columnwidth}
\caption{Where the labels go. (a)~Share of the improvement of the pooled classifier over the source-only classifier that comes from the threshold, against $K$ (median over ordered pairs without the two pairs of IEMOCAP sessions; ten iterations). Solid lines: single-modality audio and face; dashed lines: fused and text configurations; the horizontal black line marks one half. (b)~Expected residual errors after estimation from $K$ labels per emotion: the threshold offset falls as $1/(2K)$ and the direction error as $q/K$, with $q$ fitted per configuration.}
\label{fig:thresholdshare}
\end{figure}

\begin{table}[t]
\caption{Part~I quantities per configuration. Losses of the source-only classifier in points of pairwise accuracy; $q$ is the effective dimension of the direction in \eqref{eq:dircost}; the last three columns give the share of the pooled improvement that comes from the threshold (median over ordered pairs).}
\label{tab:part1}
\centering
\footnotesize
\setlength{\tabcolsep}{3.5pt}
\begin{tabular}{@{}lrrrrrr@{}}
\toprule
 & \multicolumn{2}{c}{Loss (points)} & & \multicolumn{3}{c}{Threshold share} \\
\cmidrule(lr){2-3}\cmidrule(l){5-7}
Config. & Thr. & Dir. & $q$ & $K{=}1$ & $K{=}5$ & $K{=}10$ \\
\midrule
A & \NUM{p1-e13-loss-offset-a} & \NUM{p1-e13-loss-dir-a} & \NUM{p1-qdir-a} & \NUM{p1-e29-critshare-k1-a} & \NUM{p1-e29-critshare-k5-a} & \NUM{p1-e29-critshare-k10-a} \\
A (HuBERT) & \NUM{p1-e13-loss-offset-ah} & \NUM{p1-e13-loss-dir-ah} & \NUM{p1-qdir-ah} & \NUM{p1-e29-critshare-k1-ah} & \NUM{p1-e29-critshare-k5-ah} & \NUM{p1-e29-critshare-k10-ah} \\
V & \NUM{p1-e13-loss-offset-v} & \NUM{p1-e13-loss-dir-v} & \NUM{p1-qdir-v} & \NUM{p1-e29-critshare-k1-v} & \NUM{p1-e29-critshare-k5-v} & \NUM{p1-e29-critshare-k10-v} \\
AV & \NUM{p1-e13-loss-offset-av} & \NUM{p1-e13-loss-dir-av} & \NUM{p1-qdir-av} & \NUM{p1-e29-critshare-k1-av} & \NUM{p1-e29-critshare-k5-av} & \NUM{p1-e29-critshare-k10-av} \\
T & \NUM{p1-e13-loss-offset-t} & \NUM{p1-e13-loss-dir-t} & -- & \NUM{p1-e29-critshare-k1-t} & \NUM{p1-e29-critshare-k5-t} & \NUM{p1-e29-critshare-k10-t} \\
VT & \NUM{p1-e13-loss-offset-vt} & \NUM{p1-e13-loss-dir-vt} & -- & \NUM{p1-e29-critshare-k1-vt} & \NUM{p1-e29-critshare-k5-vt} & \NUM{p1-e29-critshare-k10-vt} \\
AT & \NUM{p1-e13-loss-offset-at} & \NUM{p1-e13-loss-dir-at} & -- & \NUM{p1-e29-critshare-k1-at} & \NUM{p1-e29-critshare-k5-at} & \NUM{p1-e29-critshare-k10-at} \\
AVT & \NUM{p1-e13-loss-offset-avt} & \NUM{p1-e13-loss-dir-avt} & -- & \NUM{p1-e29-critshare-k1-avt} & \NUM{p1-e29-critshare-k5-avt} & \NUM{p1-e29-critshare-k10-avt} \\
\bottomrule
\end{tabular}
\end{table}

\subsection{How Pooling Corrects the Threshold}
\label{sec:part1-pooling}

Pooling corrects the threshold without a separate intercept for the target, and one Newton step from the source-only solution shows how (P3).
Let $N_s$ and $N_t$ be the numbers of source and target training utterances, $\kappa_s$ and $\kappa_t$ the mean curvatures of the loss at the source-only solution, $A_s$ the curvature-weighted scatter of the source features, and $\boldsymbol{\delta}$ the difference between the target mean and the curvature-weighted source mean.
One step leaves the fraction $1 - \lambda$ of the source-only offset $o_s$, with
\begin{equation}
\label{eq:lambda}
\lambda = \frac{\rho}{1 + \rho},\qquad \rho = N_t\kappa_t\left[\frac{1}{N_s\kappa_s} + \boldsymbol{\delta}^{\top}\left(A_s + \frac{I}{C}\right)^{-1}\boldsymbol{\delta}\right].
\end{equation}
The first term in the bracket is the intercept path, through which the target corrects the threshold in proportion to its share of the training weight.
The second term is a weight path: when the target is linearly separable from the source along directions in which the source varies little, the pooled classifier places weight on such a direction, and that weight acts as an intercept that only the target utterances see.
This form is the group influence of a set of points~\cite{CITE:koh2019,CITE:pregibon1981} written for added points, and its intercept path reduces to the threshold shift $m/(n+m)$ of a one-dimensional discriminant pooled from $n$ target and $m$ source points~\cite{CITE:desilva2023}.

The measured correction $\lambda = 1 - o_m/o_s$ of pooled classifiers, computed from their $d$ and $o$ on the target, is ordered as \eqref{eq:lambda} predicts, although the expression overestimates its average (Fig.~S1 of the supplementary material).
Over ordered pairs and $K \in \{1, 5, 10\}$, the rank correlation between the predicted and the measured $\lambda$ is \NUM{p1-e15-lam-rho-lo}--\NUM{p1-e15-lam-rho-hi} across A, A~(HuBERT), V, and AV (ten iterations).
The full Newton step orders the corrections with a rank correlation of \NUM{p1-e16-lam-rho-lo}--\NUM{p1-e16-lam-rho-hi}, and it also predicts the change of the direction for V and AV (\NUM{p1-e16-dir-rho-video-lo}--\NUM{p1-e16-dir-rho-video-hi}), but less well for audio (\NUM{p1-e16-dir-rho-audio-lo}--\NUM{p1-e16-dir-rho-audio-hi}), where the second-order expansion is less accurate.
The weight path explains why the distance between corpora does not predict the gain: a large distance lowers the source-only accuracy but also widens the path through which pooling corrects the threshold.
The derivation and the simulations that check \eqref{eq:lambda} are in the supplementary material.

\subsection{What a Source Is Worth}
\label{sec:part1-value}

A source contributes a direction, and its value can be expressed in target labels (P4).
Let $e$ be the efficiency of the source direction for the target, the factor by which it scales the target's best separation, $d = e\Delta$.
Under \eqref{eq:dircost}, the direction estimated from $K$ labels has a noise-to-signal ratio $q/K$, and the source direction has the ratio $(1 - e^2)/e^2$.
Combining two independent estimates with optimal weights adds their sample sizes, so starting from the source and refitting on $K$ labels behaves like training on the target alone with $K + n$ labels per emotion, where
\begin{equation}
\label{eq:value}
n = q\,\frac{e^2}{1 - e^2}.
\end{equation}
Equation~\eqref{eq:value} inverts the classical relation between the cosine of an estimated direction and the sample size~\cite{CITE:raudys1998,CITE:dobriban2018}; the inversion turns the source's direction into a count of labels, and $n$ does not depend on $K$.

We measure $n$ by locating the UAR of hypothesis transfer on the target-only learning curve and subtracting $K$ (Fig.~\ref{fig:sourcevalue}).
For a single audio source, $n$ is \NUM{p1-e32-n1src-audio-lo}--\NUM{p1-e32-n1src-audio-hi} labels per emotion and stays nearly constant from $K=5$ to $K=25$; for the face it is \NUM{p1-e32-n1src-v-k5}--\NUM{p1-e32-n1src-v-k10}.
Starting from all sources pooled, $n$ is \NUM{p1-e32-nall-a-k5} and \NUM{p1-e32-nall-a-k10} labels per emotion for WavLM audio at $K=5$ and $K=10$.
Across audio pairs, the rank correlation between the predicted and the measured $n$ is \NUM{p1-e32-rho-audio-lo}--\NUM{p1-e32-rho-audio-hi}, and the measured values are \NUM{p1-e32-ratio-audio-lo}--\NUM{p1-e32-ratio-audio-hi} times the predicted ones; for all sources pooled the ratio is \NUM{p1-e32-allsrc-ratio-lo}--\NUM{p1-e32-allsrc-ratio-hi}.

\begin{figure}[t]
\centering
\FIG{sourcevalue}{\columnwidth}
\caption{What a source is worth. Predicted value $n = q e^2/(1-e^2)$ against the value measured on the target-only learning curve, in labels per emotion, for single sources and for all sources pooled ($K=5$ and $K=10$; pairs without the two pairs of IEMOCAP sessions). The diagonal marks agreement; points beyond the range of the learning curve are shown in the top band and points with $n \le 0$ in the bottom band.}
\label{fig:sourcevalue}
\end{figure}

\subsection{Summary of Part~I}
\label{sec:part1-summary}

For single-modality audio and face features, a few target labels fix the thresholds, and centering removes most of the threshold error without any label.
The direction is the expensive part: audio needs on the order of $q$ labels per emotion to approach the target's best direction, and a source supplies about ten of them, all sources together two to three times as many.
Further progress with $K \le 5$ therefore requires more direction information than the labels and the sources provide in this form.
Part~II derives where that information can come from.
